% summary-preamble.tex
% Included in the plain-language summary build via --include-in-header.
% Intentionally omits \sectionbreak (\clearpage) — the summary is a
% shorter, reader-friendly document where forced page-breaks per section
% would be unnecessarily disruptive.

% Unicode character mappings for XeLaTeX
% The plain-language summary uses Unicode symbols directly in text mode:
%   - Proof-step marker: ✓
%   - Comparison operators: ≠ ≡ ≈ ≥ ≤
%   - Greek letter shorthand: δ
%   - Unicode superscripts for exponents: ⁰¹²³⁴⁵⁶⁷⁸⁹⁻
%     (e.g. 10⁻⁷, 10²⁹, 10³⁵ rather than LaTeX math mode)
% lmroman lacks these glyphs; newunicodechar routes them to proper LaTeX
% commands so no glyph-missing warnings are emitted.
\usepackage{amssymb}
\usepackage{newunicodechar}
% Proof marker and comparison operators
\newunicodechar{✓}{$\checkmark$}
\newunicodechar{≠}{$\neq$}
\newunicodechar{≡}{$\equiv$}
\newunicodechar{≈}{$\approx$}
\newunicodechar{≥}{$\geq$}
\newunicodechar{≤}{$\leq$}
% Greek shorthand
\newunicodechar{δ}{$\delta$}
% Unicode superscript digits and minus (used for exponents like 10⁻⁷, 10²⁹)
\newunicodechar{⁰}{$^{0}$}
\newunicodechar{¹}{$^{1}$}
\newunicodechar{²}{$^{2}$}
\newunicodechar{³}{$^{3}$}
\newunicodechar{⁴}{$^{4}$}
\newunicodechar{⁵}{$^{5}$}
\newunicodechar{⁶}{$^{6}$}
\newunicodechar{⁷}{$^{7}$}
\newunicodechar{⁸}{$^{8}$}
\newunicodechar{⁹}{$^{9}$}
\newunicodechar{⁻}{$^{-}$}
